%% related_work.tex

\section{Related Work}
\label{sec:related-work}

We position our work along three axes: prior studies of security vulnerabilities in LLM agents, the emerging literature on self-healing and autonomous error recovery in agent systems, and the long tradition of privilege management in operating systems. Our contribution sits at the intersection of these three bodies of work: we show that the \emph{recovery channel}---an artifact of self-healing design---constitutes an attack surface that prior security work has not examined, and we demonstrate that classical OS-principled privilege invariants port naturally to close it.

\subsection{Security Vulnerabilities in LLM Agents}
\label{sec:related-injection}

\textbf{Prompt injection.} The canonical vulnerability class for LLM agents is \emph{prompt injection}, in which adversary-controlled text is incorporated into the model's context and subsequently obeyed as if it were a trusted instruction. Greshake et al.\ introduced \emph{indirect prompt injection}, demonstrating that a fetched web page can embed directives that the agent executes with the privileges of its system prompt~\cite{greshake2023indirect}. Willison~\cite{willison2023prompt} catalogued the practical exploitability of injection against production agent frameworks. Liu et al.~\cite{liu2024formalizing} formalized and benchmarked prompt injection attacks and defenses, confirming that injection succeeds with high probability against virtually all production agent architectures.

\textbf{Tool manipulation.} A complementary line of work examines attacks against the \emph{tool layer} rather than the prompt layer. \cite{tool-poisoning} demonstrates that malicious MCP servers can manipulate tool descriptions to induce harmful tool calls; \cite{tool-slots} shows that adversarial tool outputs can hijack subsequent agent reasoning. These works share a common threat model: the attacker controls content that flows \emph{into} the agent---either as input, as retrieved context, or as tool return values---and the defense is therefore placed at the corresponding entry points (input sanitization~\cite{spotlighting}, instruction-hierarchy enforcement~\cite{instruction-hierarchy}, capability-based data-flow control~\cite{debenedetti2025camel}).

\textbf{Research gap.} Despite the maturity of this literature, every prior attack and defense we are aware of operates on the \emph{forward execution path}: content enters the agent, is processed, and either is obeyed or rejected at the entry boundary. We observe that this framing leaves a structural blind spot. When an agent fails---whether because of injected content or for entirely benign reasons---the failure payload (error trace, core dump, diagnostic log) re-enters the agent context via a \emph{recovery path} that has no entry-point defenses. Content that was correctly quarantined as untrusted on the forward path is re-laundered as trusted on the recovery path, because the recovery machinery was never designed to receive untrusted input. No prior work, to our knowledge, has examined the recovery path as an injection surface. This paper closes that gap.

Table~\ref{tab:related-comparison} summarizes the specific delta between our contribution and the closest prior work. The key distinctions are: (i)~we target the \emph{recovery path} (error/retry channel), not the forward execution path; (ii)~our V2 defense is \emph{structural} (privilege-downgrade invariant enforced at the framework level), not content-based (which is inherently vulnerable to paraphrasing); and (iii)~we provide a real end-to-end exploit on unmodified LangGraph, not just a conceptual argument.

\begin{table}[htbp]
\centering
\small
\caption{Comparison with prior prompt-injection and agent-security work. ``Forward path'' = defense at input/tool-entry boundary; ``Recovery path'' = defense at error/retry/repair boundary. ``Content-level'' = defense relies on detecting adversarial text; ``Structural'' = defense relies on framework-level invariants (provenance tags, privilege checks) independent of text content.}
\label{tab:related-comparison}
\setlength{\tabcolsep}{3pt}
\resizebox{\textwidth}{!}{%
\begin{tabular}{lcccc}
\toprule
\textbf{Work} & \textbf{Injection Surface} & \textbf{Defense Path} & \textbf{Defense Level} & \textbf{Real Framework Exploit?} \\
\midrule
Greshake et al.~\cite{greshake2023indirect} & Indirect (retrieved docs) & Forward & Content-level & No (conceptual) \\
Spotlighting~\cite{spotlighting} & Indirect (tool output) & Forward & Content-level & No \\
Instruction Hierarchy~\cite{instruction-hierarchy} & Direct + indirect & Forward & Content-level & No \\
CaMeL~\cite{debenedetti2025camel} & Tool invocation & Forward & Structural (perm.) & No \\
Liu et al.~\cite{liu2024formalizing} & Taxonomy (all forward) & Forward & --- & No \\
\midrule
\textbf{This work (V1)} & \textbf{Recovery (error msg)} & \textbf{Recovery} & \textbf{Content-level} & \textbf{Yes (LangGraph library-default; Reflexion commit \texttt{218cf0e})} \\
\textbf{This work (V2)} & \textbf{Recovery (core dump)} & \textbf{Recovery} & \textbf{Structural (priv.)} & \textbf{Partial (simplified AutoGen port; LangGraph author-constructed node)} \\
\bottomrule
\end{tabular}%
}
\end{table}

\subsection{Self-Healing and Autonomous Error Recovery}
\label{sec:related-selfhealing}

\textbf{Reflection and self-correction.} A substantial body of work improves LLM agent robustness by having the model reflect on its own failures and retry. Reflexion~\cite{reflexion} augments the agent with a verbal reinforcement loop that feeds a self-generated critique into the next iteration; Self-Refine~\cite{self-refine} iterates refinement with self-produced feedback; CRITIC~\cite{critic} externalizes verification to tool-augmented self-critique. These mechanisms uniformly assume that the failure payload (the critique, the error message, the verification trace) is \emph{trustworthy diagnostic signal} to be acted upon, not adversarial content to be defended against.

\textbf{Topology mutation and architectural self-repair.} Multi-agent frameworks additionally restructure the agent graph on failure. AutoGen~\cite{autogen} supports dynamic agent role reassignment; SWE-agent~\cite{swe-agent} reconfigures its tool interface when an action fails; \cite{omo} and \cite{meta-agent} describe orchestrators that swap in higher-capability models or replace a failing agent with a different role. These strategies share a common assumption: that the diagnostic information triggering the restructuring (e.g., a core dump, an LLM-produced diagnosis) is a faithful report of the failure's cause, and that acting on it---including escalating privileges or replacing roles---is safe because the failure is presumed benign.

\textbf{Implicit trust in recovery.} The robustness gains of self-healing are well-documented, but the security implications of its implicit trust assumptions have not been scrutinized. Recovery strategies systematically upgrade the trust level of the content they process: error traces are wrapped in system-level framing (``\code{[SYSTEM CRITICAL]}''); core dumps are fed to diagnostic LLMs whose outputs are treated as authoritative; role replacements suggested by the diagnostic module are executed without re-validation. Each of these is a load-bearing trust assumption that holds only if the failure payload is benign. Our work is the first to show that this assumption is violated precisely when the failure is caused by the adversary---and that the violation is exploitable with certainty.

\textbf{Concurrent and adjacent work (novelty scope).} We conducted a systematic search for prior work targeting the recovery/reflection/retry channel in LLM agents, using the keywords ``reflection injection'', ``error message injection agent'', ``retry loop prompt injection'', ``self-reflection attack'', and ``recovery channel LLM'' across arXiv (2023--2026), Google Scholar, and the major security venues. We found no prior work that (i)~formalizes a recovery-channel attack surface distinct from forward-path injection, (ii)~demonstrates end-to-end exploitability via the recovery pipeline on a production agent framework, or (iii)~proposes a privilege-downgrade invariant for recovery actions. The closest adjacent works we are aware of operate on the \emph{forward} path and are complementary to ours. Zhan et al.~\cite{zhan2025adaptive} evaluate 8 defenses against indirect prompt injection and bypass all of them with adaptive attacks (ASR~$>$~50\%); their finding that content-level forward-path defenses fall to adaptive attacks is consistent with our V1 result that content-level \emph{recovery-path} defenses degrade on stronger instruction-following models (0.07~$\rightarrow$~0.50 ASR from DeepSeek/Qwen to gpt-oss-120b, \S\ref{sec:closed-model}), but they do not consider the recovery channel. CaMeL~\cite{debenedetti2025camel} proposes a capability-based defense that separates control and data flows on the forward path; our \reauthgate{} ports the same capability-based principle to the \emph{recovery} path, suggesting the two defenses are complementary (CaMeL guards the forward execution; \reauthgate{} guards the reverse recovery). The OWASP Top 10 for LLM Applications~\cite{owasp-llm-top10} enumerates \emph{LLM01: Prompt Injection} as a forward-path risk but does not classify the recovery channel as a distinct surface. We therefore believe our contribution is novel within the LLM-agent security literature, while acknowledging that the field is moving rapidly and concurrent work may emerge between submission and publication.

\textbf{Cross-framework evidence.} To assess whether this implicit trust assumption is specific to \system{} or systemic to the LLM agent ecosystem, we inspected the recovery paths of three mainstream frameworks: AutoGen~\cite{autogen}, LangGraph~\cite{langgraph}, and SWE-agent~\cite{swe-agent}. All three exhibit the same structural blind spot, with varying severity:
\begin{itemize}[leftmargin=*]
\sloppy
    \item \textbf{AutoGen}: \code{ConversableAgent.execute\_function} formats a tool exception as \code{"Error: \{e\}"} and returns it in the same \code{function}/\code{tool} role used for successful results. A boolean \code{is\_exec\_success} signal is computed but discarded at all five call sites (assigned to \code{\_}), so the error content enters the conversation history with no provenance tag and the same trust level as a legitimate tool return.
    \item \textbf{LangGraph}: The prebuilt \code{ToolNode} can wrap tool exceptions in a \code{ToolMessage} with \code{status="error"} metadata (when \code{handle\_tool\_errors} is enabled)---the only one of the three to mark errors at the message level. However, when a tool \emph{returns} an error-like string (the common case for tools that fetch external content and encounter a processing error), the \code{ToolMessage} status is \code{"success"}---there is no content-level trust framing to distinguish error-recovery content from legitimate output. Our real end-to-end experiments (\S\ref{sec:langgraph-real}) confirm this gap is exploitable via \emph{both} forward-path and recovery-path variants: the V1 forward-path attack achieves 0.88 pooled ASR on unmodified \code{create\_react\_agent} (tool returns an error string), and the V1 recovery-path attack achieves 0.76 pooled ASR on \code{StateGraph} (tool raises an exception $\rightarrow$ recovery node wraps in system-level framing $\rightarrow$ LLM obeys). The recovery-path experiment is the primary validation of the RCI mechanism on a third-party framework.
    \item \textbf{SWE-agent}: The retry loop in \code{Agent.get\_model\_requery\_history} appends the rendered error template as \code{\{"role": "user", "message\_type": "user"\}}---indistinguishable in role and message type from a legitimate user instruction. No provenance tag distinguishes error-recovery content from user directives.
\end{itemize}
None of the three frameworks implements content-level provenance tagging (analogous to our \trustorigin{}) or a privilege-downgrade invariant on recovery actions (analogous to our \reauthgate{}). For LangGraph, we go beyond source-level observation and provide \emph{real end-to-end exploits for both vectors}: V1 achieves 0.76 pooled ASR via a recovery-path experiment on \code{StateGraph} and 0.88 via a forward-path experiment on \code{create\_react\_agent}; V2 achieves 0.54 pooled ASR on \code{StateGraph} (\S\ref{sec:langgraph-real}), with the \reauthgate{} analog reducing V2 to 0.00. For AutoGen, a simplified port reproduces the V2 vector (\S\ref{sec:second-system}). \textbf{For SWE-agent, the evidence remains at the source-code level only: we have not built a working exploit, and a real end-to-end PoC on SWE-agent is left as future work.} The absence of provenance tagging in SWE-agent's retry loop makes it a \emph{candidate} for analogous vulnerabilities, not a confirmed instance. The LangGraph and AutoGen evidence, combined with the SWE-agent source-level observation, suggests---but does not prove---that the vulnerability class is not an artifact of \system{}'s design but a gap in how the LLM agent community currently treats the recovery channel. A confirmed SWE-agent exploit would strengthen this claim.

\subsection{Privilege Management in OS and AI}
\label{sec:related-privilege}

\textbf{OS privilege invariants.} Classical operating systems enforce a small set of invariants that have proven robust over decades. The \emph{principle of least privilege}~\cite{saltzer1975protection} holds that every component should operate with the minimum privileges required for its task; the \emph{privilege downgrade invariant} holds that an exception handler executes with the privileges of the faulting context, never higher (e.g., a userspace page fault is handled in kernel mode but cannot escalate the faulting process's privileges); and \emph{provenance tagging} holds that data copied from userspace retains a \code{PAGE\_USER} marker and may not be executed as kernel code~\cite{levine2003system,bovet2005understanding}. Sandboxing~\cite{capsules,seccomp} and capability-based systems~\cite{capabilities} extend these invariants by restricting the ambient authority of a process to a declared subset.

\textbf{Privilege in AI systems.} The AI security literature has begun importing these principles. \cite{instruction-hierarchy} enforces an instruction-hierarchy that treats system prompts as higher-privilege than user content, an analog of kernel/userspace separation; CaMeL~\cite{debenedetti2025camel} extracts control and data flows from a trusted query and uses capabilities to prevent untrusted data from influencing program flow at tool-call boundaries, an analog of capability gating. These works port the \emph{forward-path} versions of the OS invariants (input validation, capability gating) to the agent setting.

\textbf{The Confused Deputy problem.} Our V2 attack vector is a textbook instance of the \emph{confused deputy} problem, originally identified by Hardy~\cite{hardy1988confused}: a privileged program (in our setting, the recovery machinery operating with role-replacement authority) is tricked by a less-privileged caller (the attacker-controlled external content embedded in the core dump) into misusing its authority on the caller's behalf. The recovery pipeline holds the privilege to swap agent roles---a capability not granted to ordinary tool calls---and, in the baseline configuration, it exercises this capability on the basis of an LLM diagnosis whose input was supplied by the attacker. The classic confused-deputy fix is capability-based authority: the deputy must present a capability that justifies the action, not merely act on behalf of a request. Our \reauthgate{} is precisely such a capability check: it demands that the proposed role replacement be in the pre-approved downgrade set of the failing role, regardless of what the diagnostic LLM suggested. The confused-deputy framing also clarifies why the V2 fix must be \emph{structural} rather than content-based: the deputy's vulnerability lies in \emph{whose authority} it is exercising, not in the textual content of the request, so no content-level defense (sanitization, framing, instruction hierarchy) can close it. To our knowledge, this is the first explicit connection between the confused-deputy problem and LLM agent recovery channels, and it positions V2 within a long-standing line of capability-based security work~\cite{capabilities,hardy1988confused}.

\textbf{Our contribution: porting the reverse-path invariants.} The observation that motivates this paper is that OS privilege invariants apply not only on the forward execution path but also---and critically---on the \emph{reverse} recovery path. Classical kernels learned long ago that exception handlers must not escalate privileges and that faulting-context data must not be re-executed as trusted code; LLM agent systems have not yet learned the corresponding lesson. Our two defenses are direct ports of these two invariants:

\begin{itemize}[leftmargin=*]
    \item \trustorigin{} Tagging is the agent analog of \code{PAGE\_USER} provenance tagging. It records, as part of the content's type, where the content came from, and refuses to execute externally-originated content under a high-trust template---just as the kernel refuses to execute a page marked \code{PAGE\_USER} in ring 0.
    \item The \reauthgate{} is the agent analog of the exception-handler privilege invariant. It enforces $P(\text{recovery action}) \leq P(\text{failed execution})$, so a recovery strategy may remediate a failure but never escalate beyond it---just as a kernel exception handler services a fault without granting the faulting process new privileges. The same principle underlies the confused-deputy fix~\cite{hardy1988confused}: the privileged deputy may not exercise authority not granted to the caller, even when the caller's request looks superficially legitimate.
\end{itemize}

Both mechanisms are lightweight (single record allocation, $O(1)$ role-set lookup; see Section~\ref{sec:discussion}), framework-agnostic in principle (they depend only on the existence of a recovery pipeline and a permission checker, not on any specific agent architecture), and---critically---\emph{local}: they are implemented at two well-defined choke points rather than scattered across the recovery strategy layer. This locality is itself an OS principle: classical kernels enforce privilege invariants at well-defined boundaries (syscall entry, exception dispatch), not inside each driver or subsystem.

In this sense, our work is not a novel cryptographic mechanism but a \emph{transfer} of proven OS invariants to a domain where they have been conspicuously absent. The empirical result---that the transfer reduces the attack surface to near-zero (V1: 0.07 residual ASR; V2: 0.00 ASR) with zero false alarms and sub-millisecond machinery overhead---suggests that the gap was one of \emph{attention}, not of mechanism: the recovery channel was overlooked, not undefendable.
